\begin{theorem}[Ramified Hasse local geometry]
\label{mod:parameters-ramifiedhasselocalgeometry}
Let $K/F$ be a ramified cyclic extension of odd prime degree $\ell$ with
prepared lower break $t$ and prepared integral generator-uniformizer
$\varpi_K$.  Put $\varpi_F=N_{K/F}(\varpi_K)$.  Suppose $d>0$,
$t=0$ or $t+1\leq d$, and
\[
  2d'+(\ell-1)t=2\ell d.
\]
For $a\in\mathcal O_F$, set
$Y_a=\varpi_K^{d'}a$.  There is a certificate with the following data.

First, the exact elementary-symmetric expansion satisfies
\[
 N_{K/F}(1+Y_a)-\bigl(1+\operatorname{Tr}_{K/F}(Y_a)+E_2(Y_a)\bigr)
 \in\mathfrak p_F^{2d+1}.
\]
Every term $E_j(Y_a)$ with $3\leq j\leq\ell$ is certified separately:
the intermediate tame terms use the totally ramified ceiling bound, the
intermediate wild terms use the strong wild symmetric bound, and the terminal
term uses the exact norm order.  Thus no higher term is silently discarded.

The certificate retains the exact trace-ideal floor
\[
 \operatorname{Tr}_{K/F}(Y_a)\in
 \mathfrak p_F^{\left\lfloor
   \frac{d'+(\ell-1)(t+1)}{\ell}\right\rfloor},
\]
and its consequences $\operatorname{Tr}(Y_a)\in\mathfrak p_F^d$ and
$E_2(Y_a)\in\mathfrak p_F^{2d}$.  If $a\in\mathfrak p_F$, the trace and
$E_2$ statements gain the powers needed for lift independence; at a positive
break the trace already lies in $\mathfrak p_F^{d+1}$.

Define, with the displayed ordered denominator,
\[
 \epsilon=
 \frac{\operatorname{Tr}_{K/F}(\varpi_K^{2d'})}{\varpi_F^{2d}}.
\]
Then $v_F(\epsilon)=0$ exactly.  The proof uses injectivity of the
above-break graded norm to rule out depth $2d+1$ for the trace numerator.
Consequently the integral representative and its residue class are
well-defined, and
\[
 \operatorname{Tr}_{K/F}(\varpi_K^{2d'})
   =\epsilon\varpi_F^{2d}.
\]
The corresponding bilinear trace specialization for $Y_aY_b$ is retained.

At a positive break the certificate also supplies the stable-to-critical
bridge which determines this residue.  For a generator $\sigma$ and
$j\in\mathbf F_p$, put
\[
  a_j=\frac{\sigma^j(\varpi_K)}{\varpi_K}-1.
\]
The exact ramification-line calculation gives
\[
  a_j\in\mathfrak p_K^t,
  \qquad
  \overline{a_j/\varpi_K^t}=j\lambda.
\]
Thus the linear term in $\prod_j(1+a_j)$ vanishes because
$\sum_{j\in\mathbf F_p}j=0$, while every nonlinear term lies in
$\mathfrak p_K^{2t}\subseteq\mathfrak p_K^{t+1}$ by
$2t\geq t+1$.  In particular, this is an explicit ramification-line
argument, and it proves
\[
  \frac{\varpi_F}{\varpi_K^p}\equiv1
  \pmod{\mathfrak p_K^{t+1}}.
\]
Together with $2d'=p(2d-t)+t$, it yields the power congruence
\[
  \varpi_K^{2d'}\equiv
  \varpi_F^{2d-t}\varpi_K^t
  \pmod{\varpi_F^{2d-t}\mathfrak p_K^{t+1}}.
\]
The trace-ideal bound then gives the residue congruence
\[
  \frac{\operatorname{Tr}_{K/F}(\varpi_K^{2d'})}{\varpi_F^{2d}}
  \equiv
  \frac{\operatorname{Tr}_{K/F}(\varpi_K^t)}{\varpi_F^t}
  \pmod{\mathfrak p_F}.
\]
The residue of the right-hand side is the extracted critical-norm linear
coefficient.  The canonical residue-degree-one equivalence also identifies
the local-geometry $\lambda$ with
\texttt{criticalNormRamificationLambda}.  Consequently the generic critical
norm polynomial gives, directly on the exported certificate,
\[
  \bar\epsilon=-\lambda^{[K:F]-1}.
\]

When $t=0$, the tame coordinate is the residue class of the exact
uniformizer-power ratio
\[
 c=\overline{\varpi_K^{d'}/\varpi_F^d}\in k_F^\times;
\]
its lift has exact order zero.  The wild coordinate is kept distinct:
for the chosen generator $\sigma$ it is
\[
 \lambda=\overline{
   \frac{\sigma(\varpi_K)-\varpi_K}{\varpi_K^{t+1}}}\in k_F^\times.
\]
The certificate proves that $\sigma\in G_t\setminus G_{t+1}$, that the
displayed numerator has exact order $t+1$, and that its normalized lift has
order zero.  Both residue classes are transported through the canonical
residue-field equivalence in the lower-to-upper direction.

This is only the local geometry used later by
\texttt{RamifiedHassePreparation}.  In particular it neither assumes nor
proves a pointwise finite Hasse-function normal form or its phase comparison;
the tame and wild finite-case residue-polynomial identities are established
when that finite comparison is constructed.
\LeanFile{LanglandsFirstMainLemma/Parameters/RamifiedHasseLocalGeometry.lean}
\lean{LanglandsFirstMainLemma.ramifiedHasseLocalGeometry}
\lean{LanglandsFirstMainLemma.ramifiedHasse_normalized_uniformizer_congr}
\lean{LanglandsFirstMainLemma.ramifiedHasse_uniformizer_power_congr}
\lean{LanglandsFirstMainLemma.ramifiedHasse_normalizedTrace_congr_critical}
\lean{LanglandsFirstMainLemma.ramifiedHasseNormalizedTraceResidue_eq_criticalNormLinearCoefficient}
\lean{LanglandsFirstMainLemma.ramifiedHasseWildLambda_eq_criticalNormRamificationLambda}
\lean{LanglandsFirstMainLemma.RamifiedHasseLocalGeometryData.wild_epsilonResidue_eq}
\leanok
\SourceText{Theorem thm:ramified-Hasse-comparison, especially equation eq:ramified-Hasse-norm-truncation-local, the depth estimates immediately following it, equation eq:ramified-Hasse-epsilon, and the subsequent tame and wild coordinate constructions.}
\uses{mod:ramification-primecyclicpreparation,mod:ramification-normbelowbreak,mod:parameters-high,mod:parameters-ramifiedhassecomparison,mod:ramification-normabovebreak,mod:ramification-symmetricbounds,mod:ramification-traceideals,mod:ramification-lowergroups}
\FormalizationContract
The exported \texttt{RamifiedHasseLocalGeometryData} supplies exactly the
geometry fields later inserted into \texttt{RamifiedHasseComparisonData}.
It uses the source-tied norm uniformizer, keeps the trace floor and all
$j\geq3$ bounds explicit, and contains no \texttt{finiteComparison} field.
At positive break its exported residue theorem factors through the generic
critical norm API: the normalized-uniformizer proof retains the actual
$j\lambda$ ramification residues and their finite-field cancellation, and
the local wild coordinate is proved compatible with the generic
ramification coefficient.
\end{theorem}
